\documentclass[12pt,a4paper]{article}
\usepackage[utf8]{inputenc}
\usepackage[T1]{fontenc}
\usepackage[italian]{babel}
\usepackage{amsmath,amssymb}
\usepackage{geometry}
\usepackage{hyperref}
\usepackage{enumitem}
\usepackage{booktabs}
\usepackage{array}
\usepackage{xcolor}
\usepackage{tcolorbox}
\usepackage{multicol}
\usepackage{tabularx}

\geometry{margin=2.2cm}
\hypersetup{colorlinks=true, linkcolor=blue!70!black}

\tcbuselibrary{skins,breakable}

\newtcolorbox{domanda}[1]{
  colback=white, colframe=blue!50!black,
  title={\textbf{#1}}, fonttitle=\small\bfseries,
  breakable, left=4pt, right=4pt, top=3pt, bottom=3pt
}
\newtcolorbox{risposta}[1]{
  colback=green!5, colframe=green!50!black,
  title={\textbf{Soluzione -- #1}}, fonttitle=\small\bfseries,
  breakable, left=4pt, right=4pt, top=3pt, bottom=3pt
}
\newtcolorbox{regola}[1]{
  colback=yellow!8, colframe=orange!60!black,
  title={\textbf{#1}}, fonttitle=\small\bfseries,
  breakable, left=4pt, right=4pt, top=3pt, bottom=3pt
}
\newtcolorbox{vfbox}[1]{
  colback=gray!5, colframe=gray!60,
  title={\textbf{#1}}, fonttitle=\small\bfseries,
  breakable, left=4pt, right=4pt, top=3pt, bottom=3pt
}
\newtcolorbox{assocbox}[1]{
  colback=white, colframe=purple!50!black,
  title={\textbf{#1}}, fonttitle=\small\bfseries,
  breakable, left=4pt, right=4pt, top=3pt, bottom=3pt
}

\newcommand{\corr}[1]{\textcolor{green!60!black}{\textbf{#1}}}
\newcommand{\V}{\textcolor{green!60!black}{\textbf{VERO}}}
\newcommand{\F}{\textcolor{red!70!black}{\textbf{FALSO}}}
\newcommand{\blank}{\underline{\hspace{4cm}}}

\setlength{\parindent}{0pt}
\setlength{\parskip}{5pt}

\title{\textbf{Fascicolo Esercizi d'Esame}\\[6pt]
{\large Sistemi Intelligenti}\\[4pt]
\normalsize Basato sulle domande reali della professoressa}
\author{}
\date{A.A. 2025--2026}

\begin{document}
\maketitle

\begin{tcolorbox}[colback=blue!5, colframe=blue!40!black]
\textbf{Come usare questo fascicolo.} Fai tutti gli esercizi di ogni sezione \emph{prima} di guardare le soluzioni. Le soluzioni di ogni sezione si trovano tutte in fondo al documento, separate, in modo da non vederle mentre lavori.
\end{tcolorbox}

\tableofcontents
\newpage

% ====================================================================
%  ██████╗ ██╗   ██╗███████╗███████╗████████╗██╗ ██████╗ ███╗   ██╗██╗
% ██╔═══██╗██║   ██║██╔════╝██╔════╝╚══██╔══╝██║██╔═══██╗████╗  ██║██║
% ██║   ██║██║   ██║█████╗  ███████╗   ██║   ██║██║   ██║██╔██╗ ██║██║
% ██║▄▄ ██║██║   ██║██╔══╝  ╚════██║   ██║   ██║██║   ██║██║╚██╗██║██║
% ╚██████╔╝╚██████╔╝███████╗███████║   ██║   ██║╚██████╔╝██║ ╚████║██║
%  ╚══▀▀═╝  ╚═════╝ ╚══════╝╚══════╝   ╚═╝   ╚═╝ ╚═════╝ ╚═╝  ╚═══╝╚═╝
% ====================================================================

\section{Domande a Scelta Multipla}

\subsection{Ricerca e algoritmi}

\begin{domanda}{D1 -- Algoritmo informato vs non informato}
Considerando un problema di ricerca in uno spazio degli stati, che cosa differenzia un algoritmo informato da uno non informato?
\begin{enumerate}[label=\alph*.]
  \item L'utilizzo di una funzione costo
  \item L'utilizzo di un modello
  \item L'utilizzo di una funzione euristica
\end{enumerate}
\end{domanda}

\begin{domanda}{D2 -- Complessità spaziale dell'Iterative Deepening}
La complessità spaziale dell'Iterative Deepening è ($br$ = branching factor, $d$ = profondità della soluzione):
\begin{enumerate}[label=\alph*.]
  \item $O(br^d)$
  \item $O(br^{d+1})$
  \item $O(br \cdot d)$
\end{enumerate}
\end{domanda}

\begin{domanda}{D3 -- Algoritmo ottimo}
Un algoritmo di ricerca in uno spazio degli stati è \textbf{ottimo} quando:
\begin{enumerate}[label=\alph*.]
  \item Restituisce sempre una soluzione nel minor numero di passi
  \item Restituisce sempre una soluzione che minimizza il valore di $h$
  \item Restituisce sempre una soluzione di costo minimo
\end{enumerate}
\end{domanda}

\begin{domanda}{D4 -- Algoritmo più efficiente in spazio}
Quale fra i seguenti è il più efficiente in termini di \textbf{spazio}?
\begin{enumerate}[label=\alph*.]
  \item UCS (Ricerca a costo uniforme)
  \item BFS (Ricerca in ampiezza)
  \item DFS (Ricerca in profondità)
\end{enumerate}
\end{domanda}

\subsection{Euristiche}

\begin{domanda}{D5 -- Euristica ammissibile}
Cosa significa che un'euristica è \textbf{ammissibile}?
\begin{enumerate}[label=\alph*.]
  \item Approssima $h^*(n)$ in maniera ottimale
  \item È sempre minore o uguale a $h^*(n)$ (non sovrastima mai)
  \item Fa stime conservative (può sottostimare)
\end{enumerate}
\end{domanda}

\begin{domanda}{D6 -- Euristica dominante}
Date due euristiche ammissibili $h_1$ e $h_2$, quando $h_1$ \textbf{domina} $h_2$?
\begin{enumerate}[label=\alph*.]
  \item Per ogni nodo $n$: $h_1(n) \geq h_2(n)$
  \item Per almeno un nodo $n$: $h_1(n) \geq h_2(n)$
  \item Per almeno un nodo $n$: $h_1(n) < h_2(n)$
\end{enumerate}
\end{domanda}

\begin{domanda}{D7 -- Euristica che non guida la ricerca}
Quale delle seguenti euristiche \textbf{non} guida l'algoritmo verso la soluzione?
\begin{enumerate}[label=\alph*.]
  \item Distanza in linea d'aria
  \item Funzione costante $h(n) = 0$
  \item Distanza di Manhattan
\end{enumerate}
\end{domanda}

\subsection{Minimax e giochi}

\begin{domanda}{D8 -- Cosa fa Minimax}
L'algoritmo Minimax\ldots
\begin{enumerate}[label=\alph*.]
  \item Costruisce un percorso che porta il giocatore Max a vincere
  \item Sceglie la prossima mossa che il giocatore Max dovrebbe eseguire
  \item Costruisce un percorso che porta il giocatore Min a perdere
\end{enumerate}
\end{domanda}

\begin{domanda}{D9 -- Come si calcolano i valori in Minimax}
I valori dei nodi di un albero costruito dall'algoritmo Minimax vengono calcolati:
\begin{enumerate}[label=\alph*.]
  \item Risalendo dalle foglie verso la radice
  \item Discendendo dalla radice verso le foglie
  \item Non si calcolano: sono dati dalla funzione di valutazione definita insieme al problema
\end{enumerate}
\end{domanda}

\begin{domanda}{D10 -- Alpha-Beta: utilità per quali nodi}
In Minimax con alpha-beta pruning, l'utilità viene calcolata:
\begin{enumerate}[label=\alph*.]
  \item Per i nodi foglia
  \item Per i nodi interni dell'albero
  \item Per la radice dell'albero
\end{enumerate}
\end{domanda}

\subsection{CSP}

\begin{domanda}{D11 -- Definizione di soluzione in un CSP}
Cos'è una \textbf{soluzione} di un CSP (problema di soddisfacimento di vincoli)?
\begin{enumerate}[label=\alph*.]
  \item Un assegnamento consistente (anche parziale)
  \item Una sequenza di passi che permette di raggiungere il goal minimizzando il costo
  \item Un assegnamento di tutte le variabili che soddisfa tutti i vincoli
\end{enumerate}
\end{domanda}

\begin{domanda}{D12 -- Algoritmo più efficiente per CSP}
Quale fra i seguenti è il \textbf{più efficiente} in termini di tempo per risolvere un CSP?
\begin{enumerate}[label=\alph*.]
  \item Ricerca con backtracking che sfrutta la commutatività
  \item Back-jumping
  \item Backward chaining
  \item Generate and test
\end{enumerate}
\end{domanda}

\begin{domanda}{D13 -- Algoritmo meno efficiente per CSP}
Quale fra i seguenti è il \textbf{meno efficiente} in termini di tempo per risolvere un CSP?
\begin{enumerate}[label=\alph*.]
  \item Ricerca con backtracking che sfrutta la commutatività
  \item Back-jumping
  \item Generate and test
\end{enumerate}
\end{domanda}

\subsection{Machine Learning e reti neurali}

\begin{domanda}{D14 -- Criterio di terminazione alberi di decisione}
Quale dei seguenti può essere usato come criterio di terminazione nella costruzione di alberi di decisione (pre-pruning)?
\begin{enumerate}[label=\alph*.]
  \item Il valore della funzione di valutazione calcolata nel nodo è maggiore di una certa soglia
  \item Il nodo presenta overfitting
  \item Il numero delle istanze associate al nodo è minore di una certa soglia
\end{enumerate}
\end{domanda}

\begin{domanda}{D15 -- Matrice di confusione: distribuzione ottimale}
Quale distribuzione dei valori in una matrice di confusione denota una \textbf{prestazione ottimale}?
\begin{enumerate}[label=\alph*.]
  \item Distribuzione uniforme del numero di istanze su tutte le celle
  \item Ogni cella contiene il numero di istanze diviso l'entropia della classe
  \item Diagonale principale che somma al numero totale delle istanze, 0 nelle altre celle
\end{enumerate}
\end{domanda}

\begin{domanda}{D16 -- Cosa calcola la matrice di confusione}
Quale delle seguenti misure viene calcolata tramite una matrice di confusione?
\begin{enumerate}[label=\alph*.]
  \item Information gain
  \item Accuratezza
  \item Entropia
\end{enumerate}
\end{domanda}

\begin{domanda}{D17 -- Cosa misura l'entropia}
Cosa misura l'entropia in un albero di decisione?
\begin{enumerate}[label=\alph*.]
  \item La presenza di più classi in un insieme di dati (grado di impurità)
  \item L'errore calcolato sull'output prodotto dalla rete neurale
  \item Un valore che permette di aggiornare i pesi di una rete neurale
\end{enumerate}
\end{domanda}

\begin{domanda}{D18 -- Cosa misura l'information gain}
Cosa misura l'information gain?
\begin{enumerate}[label=\alph*.]
  \item Il grado di purezza di un insieme di dati
  \item La riduzione di overfitting comportata dal pruning
  \item La riduzione di entropia prodotta da uno split
\end{enumerate}
\end{domanda}

\begin{domanda}{D19 -- Problema rilevante nello studio del percettrone}
Quale dei seguenti problemi è stato rilevante nello studio del percettrone?
\begin{enumerate}[label=\alph*.]
  \item Problema dell'or esclusivo (XOR)
  \item Problema della colorazione della mappa dell'Australia
  \item Anomalia di Sussman
\end{enumerate}
\end{domanda}

\begin{domanda}{D20 -- K-NN: su cosa si basa}
Sulla base di cosa l'algoritmo K Nearest Neighbour restituisce la classe di un'istanza?
\begin{enumerate}[label=\alph*.]
  \item Combinazione lineare dei pesi delle connessioni moltiplicati per gli input
  \item Soddisfacimento dell'antecedente di una regola indotta dal learning set
  \item Somiglianza dell'istanza alle $k$ istanze memorizzate più simili
\end{enumerate}
\end{domanda}

\subsection{Logica e FOL}

\begin{domanda}{D21 -- Test di Turing}
Il test di Turing dice che un computer è intelligente se:
\begin{enumerate}[label=\alph*.]
  \item Comprende le domande che gli vengono poste
  \item Imita il comportamento umano nel rispondere alle domande
  \item Risponde correttamente alle domande che gli vengono poste
  \item È in grado di manifestare emozioni
\end{enumerate}
\end{domanda}

\begin{domanda}{D22 -- Frame problem}
Per ``frame problem'' si intende:
\begin{enumerate}[label=\alph*.]
  \item Il fatto che solo alcuni concetti di un'ontologia hanno un corrispondente in un'altra
  \item L'incapacità di sapere quali proprietà di una situazione si preservano dopo l'applicazione di un'azione
  \item L'imposizione di un orizzonte temporale alla ricerca di una soluzione
\end{enumerate}
\end{domanda}

\begin{domanda}{D23 -- Unificatore nella FOL}
Qual è la definizione di \textbf{unificatore} di due formule nella logica del primo ordine?
\begin{enumerate}[label=\alph*.]
  \item Associazione fra i concetti affini di due KB differenti
  \item Sostituzione delle variabili che rende identiche le due formule
  \item Regola di inferenza che, date due clausole, restituisce il loro risolvente
\end{enumerate}
\end{domanda}

\begin{domanda}{D24 -- Cosa richiede la risoluzione}
La regola di risoluzione richiede formule espresse come:
\begin{enumerate}[label=\alph*.]
  \item Clausole di Horn
  \item Clausole qualsiasi (in CNF)
  \item Formule in forma normale disgiuntiva
\end{enumerate}
\end{domanda}

\begin{domanda}{D25 -- Path consistency}
La path consistency (consistenza di cammino):
\begin{enumerate}[label=\alph*.]
  \item Verifica la possibilità di assegnare a una variabile un valore consistente rispetto ad altre \textbf{due} variabili
  \item Verifica la possibilità di assegnare a una variabile un valore consistente rispetto ad \textbf{un'altra sola} variabile
  \item Verifica la possibilità di assegnare a una variabile un valore consistente (senza riferimento ad altre)
\end{enumerate}
\end{domanda}

\begin{domanda}{D26 -- Situation Calculus: elemento non incluso nell'azione}
Quale dei seguenti \textbf{non} fa parte della definizione di un'azione nel situation calculus?
\begin{enumerate}[label=\alph*.]
  \item Proprietà che devono valere per applicare l'azione (precondizioni)
  \item Proprietà atemporali (eternali)
  \item Proprietà che l'applicazione dell'azione modifica (effetti)
\end{enumerate}
\end{domanda}

\begin{domanda}{D27 -- RDF (Resource Description Framework)}
Considerando RDF, quale delle seguenti affermazioni è vera?
\begin{enumerate}[label=\alph*.]
  \item La conoscenza è costituita da triple $\langle$soggetto, predicato, oggetto$\rangle$
  \item La conoscenza è costituita da relazioni IS-A e MEMBER
  \item La conoscenza è costituita da clausole di Horn
\end{enumerate}
\end{domanda}

\begin{domanda}{D28 -- Tassonomie}
Quale delle seguenti affermazioni è corretta quando si parla di tassonomie?
\begin{enumerate}[label=\alph*.]
  \item L'insieme di tutte le categorie costituisce una partizione
  \item Le sottocategorie di ciascuna categoria costituiscono una partizione
  \item Tutte le categorie di una tassonomia sono disgiunte
\end{enumerate}
\end{domanda}

\begin{domanda}{D29 -- Modello nella logica proposizionale}
Quale dei seguenti potrebbe essere un modello nella logica proposizionale?
\begin{enumerate}[label=\alph*.]
  \item $X$ vale \textit{true}, $Y$ vale \textit{false}
  \item $X$ vale Giovanni, $Y$ vale Riccardo
  \item $X$ vale 4, $Y$ vale 2
\end{enumerate}
\end{domanda}

\begin{domanda}{D30 -- In quale contesto abbiamo parlato di post-pruning?}
In quale contesto abbiamo parlato di post-pruning?
\begin{enumerate}[label=\alph*.]
  \item Problemi di soddisfacimento dei vincoli
  \item Apprendimento automatico
  \item Problemi con avversario
\end{enumerate}
\end{domanda}

\newpage

% ====================================================================
\section{Esercizi Vero/Falso su Formule FOL}
% ====================================================================

\begin{regola}{Regole chiave -- leggi prima di fare gli esercizi}
\begin{itemize}[noitemsep]
  \item \textbf{Clausola:} qualsiasi disgiunzione di letterali (atomi positivi o negati).
  \item \textbf{Clausola di Horn:} clausola con \textbf{al più un} letterale positivo (non negato).
  \item \textbf{Clausola definita:} clausola di Horn con \textbf{esattamente uno} letterale positivo (ha una ``testa'').
  \item \textbf{Forward chaining / MPG applicabile:} serve una clausola definita + fatti ground che soddisfano \emph{tutte} le premesse.
  \item \textbf{Risoluzione applicabile:} servono due clausole con letterali complementari unificabili.
\end{itemize}
\end{regola}

\begin{vfbox}{Esercizio F1 -- Due formule{,} premessa incompleta}
\textbf{A)} $\neg\text{Pred1}(x) \lor \text{Pred2}(x) \lor \neg\text{Pred3}(x)$ \qquad
\textbf{B)} $\text{Pred1}(\text{ID0})$
\smallskip

Indica VERO o FALSO per ciascuna affermazione:
\begin{enumerate}[label=\arabic*.]
  \item La formula A è una clausola \hfill \underline{\hspace{2cm}}
  \item Il forward chaining è applicabile \hfill \underline{\hspace{2cm}}
  \item La formula A è una clausola di Horn \hfill \underline{\hspace{2cm}}
  \item La risoluzione è applicabile \hfill \underline{\hspace{2cm}}
  \item Il modus ponens generalizzato è applicabile \hfill \underline{\hspace{2cm}}
\end{enumerate}
\end{vfbox}

\begin{vfbox}{Esercizio F2 -- Tre formule{,} clausola di Horn definita con fatti completi}
\textbf{A)} $\neg\text{Alfa}(x) \lor \text{Beta}(x) \lor \neg\text{Gamma}(x)$ \qquad
\textbf{B)} $\text{Alfa}(\text{ID0})$ \qquad
\textbf{C)} $\text{Gamma}(\text{ID0})$
\smallskip

Indica VERO o FALSO per ciascuna affermazione:
\begin{enumerate}[label=\arabic*.]
  \item La formula A è una clausola di Horn \hfill \underline{\hspace{2cm}}
  \item Il forward chaining è applicabile \hfill \underline{\hspace{2cm}}
  \item La formula A è una clausola \hfill \underline{\hspace{2cm}}
  \item La risoluzione è applicabile \hfill \underline{\hspace{2cm}}
  \item Il modus ponens generalizzato è applicabile \hfill \underline{\hspace{2cm}}
\end{enumerate}
\end{vfbox}

\begin{vfbox}{Esercizio F3 -- Tre formule{,} clausola con DUE letterali positivi}
\textbf{A)} $\text{Alfa}(\text{ID1}) \lor \neg\text{Beta}(x) \lor \text{Gamma}(x)$ \qquad
\textbf{B)} $\text{Alfa}(\text{ID0})$ \qquad
\textbf{C)} $\text{Beta}(\text{ID1})$
\smallskip

Indica VERO o FALSO per ciascuna affermazione:
\begin{enumerate}[label=\arabic*.]
  \item La formula A è una clausola di Horn \hfill \underline{\hspace{2cm}}
  \item Il forward chaining è applicabile \hfill \underline{\hspace{2cm}}
  \item La formula A è una clausola \hfill \underline{\hspace{2cm}}
  \item Si può derivare $\text{Gamma}(\text{ID1})$ \hfill \underline{\hspace{2cm}}
  \item Il modus ponens generalizzato è applicabile \hfill \underline{\hspace{2cm}}
\end{enumerate}
\end{vfbox}

\begin{vfbox}{Esercizio F4 -- Applicare il MPG (con funzione Proprietario)}
\textbf{A)} $\text{Cane}(\text{Asso})$ \qquad
\textbf{B)} $\text{Custode}(\text{Proprietario}(z),\; z)$ \qquad
\textbf{C)} $\text{Cane}(x) \land \text{Custode}(y, x) \Rightarrow \text{Responsabile}(y)$
\smallskip

Indica VERO o FALSO per ciascuna affermazione:
\begin{enumerate}[label=\arabic*.]
  \item Il MPG non è applicabile \hfill \underline{\hspace{2cm}}
  \item Sostituzione corretta: $\{x/\text{Asso},\; z/z,\; y/\text{Proprietario}(\text{Asso})\}$ \hfill \underline{\hspace{2cm}}
  \item Risultato: $\text{Cane}(\text{Asso}) \land \text{Custode}(y, \text{Asso})$ \hfill \underline{\hspace{2cm}}
  \item Risultato: $\text{Responsabile}(\text{Proprietario}(\text{Asso}))$ \hfill \underline{\hspace{2cm}}
  \item Sostituzione corretta: $\{x/\text{Asso},\; z/\text{Asso},\; y/\text{Proprietario}(\text{Asso})\}$ \hfill \underline{\hspace{2cm}}
\end{enumerate}
\end{vfbox}

\begin{vfbox}{Esercizio F5 -- Applicare la Risoluzione FOL}
\textbf{A)} $\neg\text{Ottone}(x) \lor \text{Percussione}(x) \lor \neg\text{Arco}(x)$ \qquad
\textbf{B)} $\text{Ottone}(\text{ID0})$
\smallskip

Indica VERO o FALSO per ciascuna affermazione:
\begin{enumerate}[label=\arabic*.]
  \item L'unificatore è $\{x/\text{ID0}\}$ \hfill \underline{\hspace{2cm}}
  \item La risoluzione non è applicabile \hfill \underline{\hspace{2cm}}
  \item Risolvente: $\text{Ottone}(\text{ID0})$ \hfill \underline{\hspace{2cm}}
  \item Risolvente: $\text{Percussione}(\text{ID0}) \lor \neg\text{Arco}(\text{ID0})$ \hfill \underline{\hspace{2cm}}
  \item Risolvente: $\text{Percussione}(x) \lor \neg\text{Arco}(x)$ \hfill \underline{\hspace{2cm}}
\end{enumerate}
\end{vfbox}

\newpage

% ====================================================================
\section{Esercizi di Associazione}
% ====================================================================

\begin{assocbox}{Associazione A1 -- Meccanismi logici e tipi di ragionamento}
Collega ogni meccanismo al tipo di ragionamento che supporta:

\begin{center}
\renewcommand{\arraystretch}{1.5}
\begin{tabular}{lll}
\textbf{Meccanismo} & & \textbf{Tipo di ragionamento} \\
\hline
Model checking & $\longrightarrow$ & \underline{\hspace{5cm}} \\
Modus Ponens Generalizzato & $\longrightarrow$ & \underline{\hspace{5cm}} \\
Proposizionalizzazione & $\longrightarrow$ & \underline{\hspace{5cm}} \\
Modus Ponens (proposizionale) & $\longrightarrow$ & \underline{\hspace{5cm}} \\
Binary Resolution & $\longrightarrow$ & \underline{\hspace{5cm}} \\
\end{tabular}
\end{center}

\smallskip
\textit{Opzioni:} conseguenza logica $\cdot$ forward chaining FOL $\cdot$ traduzione KB FOL in logica proposizionale $\cdot$ forward chaining proposizionale $\cdot$ ragionamento per refutazione in FOL
\end{assocbox}

\begin{assocbox}{Associazione A2 -- Algoritmi e compiti}
Collega ogni algoritmo al compito che svolge:

\begin{center}
\renewcommand{\arraystretch}{1.5}
\begin{tabular}{lll}
\textbf{Algoritmo} & & \textbf{Compito} \\
\hline
Algoritmo di Hunt & $\longrightarrow$ & \underline{\hspace{5cm}} \\
K-NN & $\longrightarrow$ & \underline{\hspace{5cm}} \\
Back-jumping & $\longrightarrow$ & \underline{\hspace{5cm}} \\
RBFS & $\longrightarrow$ & \underline{\hspace{5cm}} \\
Alpha-beta pruning & $\longrightarrow$ & \underline{\hspace{5cm}} \\
Minimax & $\longrightarrow$ & \underline{\hspace{5cm}} \\
Forward chaining & $\longrightarrow$ & \underline{\hspace{5cm}} \\
Backward chaining & $\longrightarrow$ & \underline{\hspace{5cm}} \\
\end{tabular}
\end{center}

\smallskip
\textit{Opzioni:} costruzione alberi di decisione $\cdot$ classificazione senza modello esplicito $\cdot$ risoluzione CSP con conflict set $\cdot$ ricerca in spazio degli stati con memoria limitata $\cdot$ riduzione spazio di ricerca in giochi $\cdot$ suggerimento prossima mossa con avversario $\cdot$ inferenza da KB in clausole di Horn $\cdot$ verifica se un goal è derivabile da KB e fatti
\end{assocbox}

\begin{assocbox}{Associazione A3 -- Euristiche e proprietà}
Collega ogni euristica/concetto alla proprietà corrispondente:

\begin{center}
\renewcommand{\arraystretch}{1.5}
\begin{tabular}{lll}
\textbf{Euristica / Concetto} & & \textbf{Proprietà} \\
\hline
Euristica dominante & $\longrightarrow$ & \underline{\hspace{5cm}} \\
Euristica ammissibile & $\longrightarrow$ & \underline{\hspace{5cm}} \\
Euristica MRV & $\longrightarrow$ & \underline{\hspace{5cm}} \\
Euristica di grado & $\longrightarrow$ & \underline{\hspace{5cm}} \\
Euristica monotona (consistente) & $\longrightarrow$ & \underline{\hspace{5cm}} \\
\end{tabular}
\end{center}

\smallskip
\textit{Opzioni:} approssima meglio $h^*(n)$ $\cdot$ è ottimistica (non sovrastima) $\cdot$ facilita il fallimento precoce $\cdot$ identifica la variabile coinvolta in più vincoli $\cdot$ garantisce ottimalità su grafi con lista chiusa
\end{assocbox}

\begin{assocbox}{Associazione A4 -- Ontologia ambito medico}
Effettua le corrette associazioni (X è una variabile):

\begin{center}
\renewcommand{\arraystretch}{1.5}
\begin{tabular}{lll}
\textbf{Espressione} & & \textbf{Tipo di relazione} \\
\hline
\texttt{Capsomero part-of Virus} & $\longrightarrow$ & \underline{\hspace{4cm}} \\
\texttt{Virologo isa Microbiologo} & $\longrightarrow$ & \underline{\hspace{4cm}} \\
\texttt{HIV member Virus} & $\longrightarrow$ & \underline{\hspace{4cm}} \\
\texttt{member(X, positivo) $\Rightarrow$ quarantenato(X)} & $\longrightarrow$ & \underline{\hspace{4cm}} \\
\texttt{\{virus, batterio, fungo, protozoo\}} rispetto a Microrganismo & $\longrightarrow$ & \underline{\hspace{4cm}} \\
\end{tabular}
\end{center}

\smallskip
\textit{Opzioni:} relazione di composizione $\cdot$ relazione di sottoclasse $\cdot$ relazione di istanza $\cdot$ proprietà di classe $\cdot$ partizione
\end{assocbox}

\begin{assocbox}{Associazione A5 -- Ontologia ambito astronomico}
Effettua le corrette associazioni (X è una variabile):

\begin{center}
\renewcommand{\arraystretch}{1.5}
\begin{tabular}{lll}
\textbf{Espressione} & & \textbf{Tipo di relazione} \\
\hline
\texttt{Mercurio member Pianeta} & $\longrightarrow$ & \underline{\hspace{4cm}} \\
\texttt{member(X, pianeta) $\Rightarrow$ sferoidale(X)} & $\longrightarrow$ & \underline{\hspace{4cm}} \\
\texttt{Astrometria isa Astronomia} & $\longrightarrow$ & \underline{\hspace{4cm}} \\
\texttt{Magnetosfera part-of Pianeta} & $\longrightarrow$ & \underline{\hspace{4cm}} \\
\texttt{\{solare, extrasolare, interstellare\}} rispetto a Pianeta & $\longrightarrow$ & \underline{\hspace{4cm}} \\
\end{tabular}
\end{center}

\smallskip
\textit{Opzioni:} relazione di composizione $\cdot$ relazione di sottoclasse $\cdot$ relazione di istanza $\cdot$ proprietà di classe $\cdot$ partizione
\end{assocbox}

\newpage

% ====================================================================
\section{Domande Aperte}
% ====================================================================

\begin{domanda}{O1 -- Tautologie{,} contraddizioni{,} soddisfacibilità}
Definisci i tre concetti e fornisci un esempio per ciascuno.

\vspace{3cm}
\end{domanda}

\begin{domanda}{O2 -- Funzioni di Skolem e termine ground}
Spiega cosa sono e come si usano nella conversione in CNF nella FOL.

\vspace{3cm}
\end{domanda}

\begin{domanda}{O3 -- Forward chaining: algoritmo e quando è applicabile}
Descrivi l'algoritmo e le condizioni necessarie per applicarlo.

\vspace{3cm}
\end{domanda}

\begin{domanda}{O4 -- Definizione di gioco a somma zero e funzionamento di Minimax}
\vspace{4cm}
\end{domanda}

\begin{domanda}{O5 -- Clausole di Horn: definizione{,} tipi e perché sono efficienti}
\vspace{4cm}
\end{domanda}

\begin{domanda}{O6 -- Overfitting: definizione{,} cause e soluzioni per gli alberi di decisione}
\vspace{4cm}
\end{domanda}

\begin{domanda}{O7 -- Back-jumping e conflict set: come funziona e perché è meglio del backtracking cronologico}
\vspace{4cm}
\end{domanda}

\begin{domanda}{O8 -- Perceptron: architettura{,} apprendimento{,} limiti}
\vspace{4cm}
\end{domanda}

% ====================================================================
% ╔═══════════════════════════════════════════════╗
% ║            SOLUZIONI -- LEGGILE DOPO          ║
% ╚═══════════════════════════════════════════════╝
% ====================================================================

\newpage
\begin{tcolorbox}[colback=red!8, colframe=red!60!black, fonttitle=\Large\bfseries,
  title=SOLUZIONI -- Apri questa sezione solo dopo aver completato tutti gli esercizi]
\end{tcolorbox}

% ====================================================================
\section*{Soluzioni Sezione 1 -- Domande a Scelta Multipla}
\addcontentsline{toc}{section}{Soluzioni}
% ====================================================================

\begin{risposta}{D1}
\textbf{Risposta: c.} Un algoritmo \emph{informato} usa una funzione euristica $h(n)$ che stima il costo residuo verso il goal. La funzione costo $g(n)$ è usata da entrambi i tipi.
\end{risposta}

\begin{risposta}{D2}
\textbf{Risposta: c} -- $O(br \cdot d)$. IDDFS mantiene in memoria solo il cammino corrente, come DFS. Il tempo è $O(b^d)$ come BFS, ma lo spazio è lineare.
\end{risposta}

\begin{risposta}{D3}
\textbf{Risposta: c.} Ottimalità = trovare sempre la \textbf{soluzione di costo minimo} (somma dei costi degli archi), non il numero minimo di passi.
\end{risposta}

\begin{risposta}{D4}
\textbf{Risposta: c (DFS).} DFS usa $O(b \cdot m)$: mantiene solo il cammino corrente + i fratelli non espansi. BFS e UCS usano $O(b^d)$.
\end{risposta}

\begin{risposta}{D5}
\textbf{Risposta: b.} Ammissibile $\Leftrightarrow$ $h(n) \leq h^*(n)$ per ogni $n$. Non sovrastima mai il costo reale. Questo garantisce l'ottimalità di A*.
\end{risposta}

\begin{risposta}{D6}
\textbf{Risposta: a.} $h_1$ domina $h_2$ se $h_1(n) \geq h_2(n)$ per \textbf{tutti} i nodi. Con $h_1$ dominante, A* espande meno nodi (stime più alte = più informative, ma ancora ammissibili).
\end{risposta}

\begin{risposta}{D7}
\textbf{Risposta: b.} Con $h(n) = 0$, A* si riduce a UCS: esplora per costo $g$ crescente senza alcuna guida verso il goal. È ammissibile ma non direzionale.
\end{risposta}

\begin{risposta}{D8}
\textbf{Risposta: b.} Minimax \emph{sceglie la prossima mossa} ottimale per Max, assumendo che Min giochi perfettamente. Non garantisce la vittoria -- suggerisce la mossa migliore \emph{qui e ora}.
\end{risposta}

\begin{risposta}{D9}
\textbf{Risposta: a.} I valori si assegnano prima alle foglie (stati terminali) e si propagano \textbf{risalendo}: nodi MAX = max dei figli, nodi MIN = min dei figli.
\end{risposta}

\begin{risposta}{D10}
\textbf{Risposta: a.} La funzione Utility si calcola solo nei nodi \textbf{foglia} (stati terminali). I valori dei nodi interni sono propagati dal basso, non calcolati direttamente.
\end{risposta}

\begin{risposta}{D11}
\textbf{Risposta: c.} La soluzione deve essere \textbf{completa} (tutte le variabili valorizzate) \emph{e} \textbf{consistente} (nessun vincolo violato). Un assegnamento solo parzialmente consistente non è una soluzione.
\end{risposta}

\begin{risposta}{D12}
\textbf{Risposta: a.} Il backtracking con commutatività è il più efficiente tra quelli elencati. Il backward chaining non è per CSP. Il generate and test è il peggiore.
\end{risposta}

\begin{risposta}{D13}
\textbf{Risposta: c.} Generate and test genera \emph{tutte} le assegnazioni complete possibili ($O(d^n)$) poi verifica i vincoli: nessuna potatura.
\end{risposta}

\begin{risposta}{D14}
\textbf{Risposta: c.} Numero di istanze sotto soglia è un criterio valido di pre-pruning. La ``funzione di valutazione'' è dei giochi. Rilevare l'overfitting a priori non è praticabile.
\end{risposta}

\begin{risposta}{D15}
\textbf{Risposta: c.} Diagonale principale = solo classificazioni corrette (TP e TN). Se la diagonale somma al totale e le altre celle sono 0 $\Rightarrow$ accuratezza 100\%.
\end{risposta}

\begin{risposta}{D16}
\textbf{Risposta: b.} Dalla matrice di confusione si calcola: accuratezza, error rate, precisione, recall. Information gain ed entropia sono per gli alberi di decisione (split).
\end{risposta}

\begin{risposta}{D17}
\textbf{Risposta: a.} $H(t) = -\sum_i p(i|t)\log_2 p(i|t)$. $H=0$ se puro, massima se classi equiprobabili. Misura il \textbf{grado di impurità} (disordine) di un nodo.
\end{risposta}

\begin{risposta}{D18}
\textbf{Risposta: c.} $\text{Gain} = H(\text{padre}) - \sum_v \frac{|D_v|}{|D_t|} H(D_v)$. Misura quanto lo split \emph{riduce} l'entropia. Si sceglie l'attributo con gain massimo.
\end{risposta}

\begin{risposta}{D19}
\textbf{Risposta: a.} Minsky e Papert (1969) dimostrarono che il percettrone singolo non può risolvere XOR (non linearmente separabile). Questo provocò il primo AI winter. Serve MLP.
\end{risposta}

\begin{risposta}{D20}
\textbf{Risposta: c.} K-NN è lazy learning: memorizza tutto il training set e classifica in base alla similarità con le $k$ istanze più vicine. Non costruisce un modello esplicito.
\end{risposta}

\begin{risposta}{D21}
\textbf{Risposta: b.} Il test è comportamentale: la macchina deve \emph{imitare} il comportamento umano abbastanza bene da ingannare il giudice. Non si richiede comprensione vera (IA forte).
\end{risposta}

\begin{risposta}{D22}
\textbf{Risposta: b.} Il frame problem è la difficoltà di rappresentare ciò che \emph{non cambia} dopo un'azione. Soluzione: assioma di stato successore nel situation calculus.
\end{risposta}

\begin{risposta}{D23}
\textbf{Risposta: b.} L'unificatore (MGU) è una sostituzione $\theta$ che rende due espressioni sintatticamente identiche: $\alpha\theta = \beta\theta$.
\end{risposta}

\begin{risposta}{D24}
\textbf{Risposta: b.} La risoluzione si applica a qualsiasi coppia di \textbf{clausole} (qualsiasi, non solo di Horn) con letterali complementari. Il MPG invece richiede clausole di Horn.
\end{risposta}

\begin{risposta}{D25}
\textbf{Risposta: a.} Path consistency (3-consistenza): per ogni assegnamento consistente di $X_i$ e $X_j$, esiste un valore per $X_m$ compatibile con entrambe. Coinvolge 3 variabili.
\end{risposta}

\begin{risposta}{D26}
\textbf{Risposta: b.} Le proprietà atemporali (eternali) non cambiano con le azioni -- non fanno parte della definizione di azione. Un'azione è descritta da precondizioni + effetti.
\end{risposta}

\begin{risposta}{D27}
\textbf{Risposta: a.} RDF rappresenta la conoscenza come grafo di triple $\langle$soggetto, predicato, oggetto$\rangle$. Ogni tripla è una asserzione.
\end{risposta}

\begin{risposta}{D28}
\textbf{Risposta: b.} Le sottocategorie di \emph{ciascuna} categoria formano una partizione (disgiunte + esaustive rispetto al padre). L'insieme di tutte le categorie non è necessariamente una partizione.
\end{risposta}

\begin{risposta}{D29}
\textbf{Risposta: a.} Nella logica proposizionale un modello assegna valori \textit{true}/\textit{false} ai simboli proposizionali. Valori oggetto (``Giovanni'', ``4'') appartengono alla FOL.
\end{risposta}

\begin{risposta}{D30}
\textbf{Risposta: b.} Il post-pruning si applica agli \textbf{alberi di decisione}: si costruisce l'albero completo poi si potano i rami ridondanti per ridurre l'overfitting.
\end{risposta}

% ====================================================================
\section*{Soluzioni Sezione 2 -- Vero/Falso FOL}
% ====================================================================

\begin{risposta}{F1}
\begin{enumerate}[noitemsep, label=\arabic*.]
  \item \V\ -- È una disgiunzione di letterali.
  \item \F\ -- Forward chaining richiede che \textbf{tutte} le premesse siano soddisfatte. A come regola: $\text{Pred1}(x) \land \text{Pred3}(x) \Rightarrow \text{Pred2}(x)$. Abbiamo B = $\text{Pred1}(\text{ID0})$ ma non $\text{Pred3}(\text{ID0})$ $\Rightarrow$ non applicabile.
  \item \V\ -- Un solo letterale positivo: \texttt{Pred2(x)}.
  \item \V\ -- A contiene $\neg\text{Pred1}(x)$, B contiene $\text{Pred1}(\text{ID0})$: complementari con $\{x/\text{ID0}\}$. Risolvente: $\text{Pred2}(\text{ID0}) \lor \neg\text{Pred3}(\text{ID0})$.
  \item \F\ -- Stessa ragione del punto 2: manca Pred3.
\end{enumerate}
\end{risposta}

\begin{risposta}{F2}
\begin{enumerate}[noitemsep, label=\arabic*.]
  \item \V\ -- Un solo letterale positivo: \texttt{Beta(x)}.
  \item \V\ -- Regola: $\text{Alfa}(x) \land \text{Gamma}(x) \Rightarrow \text{Beta}(x)$. B dà $\text{Alfa}(\text{ID0})$, C dà $\text{Gamma}(\text{ID0})$: entrambe le premesse soddisfatte con $\{x/\text{ID0}\}$. Si deriva $\text{Beta}(\text{ID0})$.
  \item \V\ -- Disgiunzione di letterali.
  \item \V\ -- A e B hanno letterali complementari $\neg\text{Alfa}(x)$ / $\text{Alfa}(\text{ID0})$ con $\{x/\text{ID0}\}$.
  \item \V\ -- Tutte le premesse soddisfatte (vedi punto 2).
\end{enumerate}
\end{risposta}

\begin{risposta}{F3}
\begin{enumerate}[noitemsep, label=\arabic*.]
  \item \F\ -- A ha \textbf{due} letterali positivi: \texttt{Alfa(ID1)} e \texttt{Gamma(x)}. Non è di Horn.
  \item \F\ -- Forward chaining e MPG richiedono clausole di Horn definite.
  \item \V\ -- È una disgiunzione di letterali.
  \item \F\ -- Non è derivabile con le regole disponibili: A non è una clausola definita.
  \item \F\ -- Come punto 2.
\end{enumerate}
\end{risposta}

\begin{risposta}{F4}
Regola C: $\text{Cane}(x) \land \text{Custode}(y, x) \Rightarrow \text{Responsabile}(y)$. Unifico le premesse:
$\text{Cane}(x)$ con A $\Rightarrow$ $\{x/\text{Asso}\}$.
$\text{Custode}(y, \text{Asso})$ con B = $\text{Custode}(\text{Proprietario}(z), z)$ $\Rightarrow$ $\{y/\text{Proprietario}(z), z/\text{Asso}\}$.
Sostituzione completa: $\{x/\text{Asso}, z/\text{Asso}, y/\text{Proprietario}(\text{Asso})\}$.
\begin{enumerate}[noitemsep, label=\arabic*.]
  \item \F\ -- Il MPG è applicabile.
  \item \F\ -- La $z$ deve essere istanziata a \texttt{Asso}, non rimane variabile libera.
  \item \F\ -- Il risultato del MPG è la conclusione istanziata, non un fatto già noto.
  \item \V\ -- Risultato corretto: $\text{Responsabile}(\text{Proprietario}(\text{Asso}))$.
  \item \V\ -- Sostituzione corretta.
\end{enumerate}
\end{risposta}

\begin{risposta}{F5}
A contiene $\neg\text{Ottone}(x)$; B contiene $\text{Ottone}(\text{ID0})$: complementari con $\{x/\text{ID0}\}$.
\begin{enumerate}[noitemsep, label=\arabic*.]
  \item \V\ -- L'unificatore è $\{x/\text{ID0}\}$ (la variabile è $x$).
  \item \F\ -- La risoluzione è applicabile (ci sono letterali complementari).
  \item \F\ -- Il risolvente non è uno dei fatti originali; si eliminano i letterali complementari.
  \item \V\ -- Risolvente corretto: $\text{Percussione}(\text{ID0}) \lor \neg\text{Arco}(\text{ID0})$ (le restanti parti di A, istanziate).
  \item \F\ -- Il risolvente è istanziato ($x$ viene sostituita da \texttt{ID0}), non rimane con variabili libere.
\end{enumerate}
\end{risposta}

% ====================================================================
\section*{Soluzioni Sezione 3 -- Associazioni}
% ====================================================================

\begin{risposta}{A1 -- Meccanismi logici}
\begin{itemize}[noitemsep]
  \item Model checking $\to$ \textbf{conseguenza logica}
  \item MPG $\to$ \textbf{forward chaining FOL}
  \item Proposizionalizzazione $\to$ \textbf{traduzione KB FOL in logica proposizionale}
  \item Modus Ponens $\to$ \textbf{forward chaining proposizionale}
  \item Binary Resolution $\to$ \textbf{ragionamento per refutazione in FOL}
\end{itemize}
\end{risposta}

\begin{risposta}{A2 -- Algoritmi e compiti}
\begin{itemize}[noitemsep]
  \item Algoritmo di Hunt $\to$ \textbf{costruzione alberi di decisione}
  \item K-NN $\to$ \textbf{classificazione senza modello esplicito (lazy learning)}
  \item Back-jumping $\to$ \textbf{risoluzione CSP con conflict set}
  \item RBFS $\to$ \textbf{ricerca in spazio degli stati con memoria limitata}
  \item Alpha-beta pruning $\to$ \textbf{riduzione spazio di ricerca in giochi}
  \item Minimax $\to$ \textbf{suggerimento prossima mossa con avversario}
  \item Forward chaining $\to$ \textbf{inferenza da KB in clausole di Horn}
  \item Backward chaining $\to$ \textbf{verifica se un goal è derivabile da KB e fatti}
\end{itemize}
\end{risposta}

\begin{risposta}{A3 -- Euristiche e proprietà}
\begin{itemize}[noitemsep]
  \item Dominante $\to$ \textbf{approssima meglio $h^*(n)$} (stima più alta, ancora ammissibile)
  \item Ammissibile $\to$ \textbf{è ottimistica} (non sovrastima mai)
  \item MRV $\to$ \textbf{facilita il fallimento precoce} (fail-first)
  \item Di grado $\to$ \textbf{identifica la variabile coinvolta in più vincoli}
  \item Monotona $\to$ \textbf{garantisce ottimalità su grafi} con lista chiusa
\end{itemize}
\end{risposta}

\begin{risposta}{A4 -- Ontologia medica}
\begin{itemize}[noitemsep]
  \item \texttt{Capsomero part-of Virus} $\to$ \textbf{relazione di composizione}
  \item \texttt{Virologo isa Microbiologo} $\to$ \textbf{relazione di sottoclasse}
  \item \texttt{HIV member Virus} $\to$ \textbf{relazione di istanza}
  \item \texttt{member(X, positivo) $\Rightarrow$ quarantenato(X)} $\to$ \textbf{proprietà di classe}
  \item \texttt{\{virus, batterio, fungo, protozoo\}} $\to$ \textbf{partizione} rispetto a Microrganismo
\end{itemize}
\end{risposta}

\begin{risposta}{A5 -- Ontologia astronomica}
\begin{itemize}[noitemsep]
  \item \texttt{Mercurio member Pianeta} $\to$ \textbf{relazione di istanza}
  \item \texttt{member(X, pianeta) $\Rightarrow$ sferoidale(X)} $\to$ \textbf{proprietà di classe}
  \item \texttt{Astrometria isa Astronomia} $\to$ \textbf{relazione di sottoclasse}
  \item \texttt{Magnetosfera part-of Pianeta} $\to$ \textbf{relazione di composizione}
  \item \texttt{\{solare, extrasolare, interstellare\}} $\to$ \textbf{partizione} rispetto a Pianeta
\end{itemize}
\end{risposta}

% ====================================================================
\section*{Soluzioni Sezione 4 -- Domande Aperte}
% ====================================================================

\begin{risposta}{O1 -- Tautologie{,} contraddizioni{,} soddisfacibilità}
\textbf{Tautologia:} vera in tutti i modelli. Es: $A \lor \neg A$. \\
\textbf{Contraddizione:} falsa in tutti i modelli. Es: $A \land \neg A$. \\
\textbf{Soddisfacibile:} vera in almeno un modello. Es: $A \land B$ (vera se entrambi true). \\
Ogni tautologia è soddisfacibile. Una contraddizione non è soddisfacibile. SAT è NP-completo.
\end{risposta}

\begin{risposta}{O2 -- Funzioni di Skolem e termine ground}
\textbf{Funzione di Skolem:} nella conversione in CNF, si elimina $\exists y$ sostituendo $y$ con una funzione $f(x_1,\ldots,x_k)$ dove $x_1,\ldots,x_k$ sono le variabili universali che contengono l'$\exists$. Se l'$\exists$ non è dentro un $\forall$, si usa una costante di Skolem $k$. \\
\textbf{Termine ground:} termine senza variabili (completamente istanziato). Es: \texttt{Padre(Giovanni)}, \texttt{Asso}. Contrapposto a \texttt{Padre(x)} (con variabile).
\end{risposta}

\begin{risposta}{O3 -- Forward chaining}
\textbf{Algoritmo:} parte dai fatti noti nella KB. Cerca clausole definite (Horn con testa) le cui premesse siano tutte soddisfatte dai fatti. Aggiunge la conclusione come nuovo fatto. Ripete fino al goal o al punto fisso (nessun nuovo fatto). \\
\textbf{Condizioni:} serve almeno una clausola definita e fatti che unifichino con \emph{tutte} le sue premesse. \\
\textbf{Proprietà:} corretto e completo per clausole definite. Diretto ai dati (data-driven).
\end{risposta}

\begin{risposta}{O4 -- Gioco a somma zero e Minimax}
\textbf{Somma zero:} ciò che guadagna MAX è perso da MIN. Informazione perfetta (tutto visibile). \\
\textbf{Minimax:} esplora l'albero in profondità. Foglie = valori di utilità. Si propaga verso l'alto: MAX prende il massimo dei figli, MIN prende il minimo. La radice ha il valore della mossa ottimale per MAX assumendo Min perfetto. Complessità: $O(b^m)$ tempo, $O(bm)$ spazio.
\end{risposta}

\begin{risposta}{O5 -- Clausole di Horn}
\textbf{Clausola di Horn:} al più un letterale positivo. Forma: $\neg P_1 \lor \cdots \lor \neg P_k \lor Q$ $\equiv$ $P_1 \land \cdots \land P_k \Rightarrow Q$. \\
\textbf{Tipi:} clausola definita (esattamente 1 positivo -- ha testa); clausola obiettivo (0 positivi); fatto (1 positivo, 0 premesse). \\
\textbf{Efficienza:} l'inferenza su clausole di Horn è lineare (polinomiale); molto più efficiente della risoluzione generale su CNF arbitraria.
\end{risposta}

\begin{risposta}{O6 -- Overfitting}
\textbf{Definizione:} il modello si adatta troppo al training set (incluso rumore) e generalizza male. \\
\textbf{Cause:} modello troppo complesso, training set troppo piccolo, presenza di rumore. \\
\textbf{Soluzioni per alberi di decisione:}
\begin{itemize}[noitemsep]
  \item \textbf{Pre-pruning:} interrompere se il nodo ha poche istanze (sotto soglia) o gain troppo basso.
  \item \textbf{Post-pruning:} costruire l'albero completo poi sostituire sottoalberi con foglie (classe più frequente).
\end{itemize}
Rasoio di Occam: a parità di errore, preferire il modello più semplice.
\end{risposta}

\begin{risposta}{O7 -- Back-jumping e conflict set}
\textbf{Problema del backtracking cronologico:} torna sempre alla variabile precedente anche se non è la causa del fallimento. \\
\textbf{Back-jumping:} quando $X$ fallisce, si costruisce il \textbf{conflict set} $CS(X)$ = variabili già assegnate che causano l'incompatibilità. L'algoritmo salta direttamente all'ultima variabile in $CS(X)$, saltando variabili irrilevanti. \\
\textbf{Implementazione:} ogni livello riceve il CS dai figli falliti, aggiunge la propria variabile e lo passa al padre.
\end{risposta}

\begin{risposta}{O8 -- Perceptron}
\textbf{Architettura:} singolo neurone. Calcola $y = f(\mathbf{w}\cdot\mathbf{x} + b)$. I pesi definiscono un iperpiano: istanze sopra $\to$ classe 1, sotto $\to$ classe 0. \\
\textbf{Apprendimento:} regola delta -- $\Delta w_j = \eta(d-o)x_j$. Aggiornamento passo per passo su ogni istanza. Converge se il problema è linearmente separabile (Teorema di Novikoff). \\
\textbf{Limiti:} non risolve XOR (non linearmente separabile). Serve MLP (con hidden layer) per funzioni non lineari.
\end{risposta}

\end{document}
